\sectionkicker{Theme synthesis}
\subsection*{横断比較 — Open Weightとmanaged lifecycleを分けて読む}
\addcontentsline{toc}{subsection}{横断比較 — Open Weightとmanaged lifecycleを分けて読む}
gpt-oss、Qwen、DeepSeek、Kimi、Mistral、GLMは一つのランキングへ並べるより、weights/license、self-host、managed API、regional availability、coding/tool surfaceという採用条件の違いとして比較した方が2025年後半の変化を捉えやすい。
\medskip
\begin{center}
\small
\renewcommand{\arraystretch}{1.16}
\begin{tabularx}{\linewidth}{@{}p{0.20\linewidth}X@{}}
\toprule
観察軸 & 一次資料・論文から読めること \\
\midrule
open-weight deployment & gpt-ossはopen-weight reasoning modelをlocal/self-hosted deploymentの選択肢として提示した。ただしweightsの公開はtraining dataや学習過程の全面開示を意味しない。 \cite{src-1c46fdf17c79,src-12865e8612c7} \\
catalog・regional availability & Alibaba Model Studioはmodel release、snapshot、region別availability、後日のglobal availabilityを同じliving catalogに持つため、catalog上の日付をupstream model releaseへ短絡させない。 \cite{src-eac8dfbeb78c} \\
API service transition & DeepSeekの更新はstable/experimental endpointやservice transitionを含む。API上の切替日はcheckpoint publicationやtraining completionの代替日付ではない。 \cite{src-e547d0d8cbe3} \\
coding/product surface & Kimi Codeの更新はcoding agent/productのliving historyとして、Mistralのnewsは公開chronologyとして有用だが、両者ともindexから得られる範囲とitem-levelの主張を分けて扱う。 \cite{src-9191d43918fc,src-d8b7c1c36830} \\
release scopeとvendor claim & GLM/AutoGLM系列はrelease identity、日付、feature scopeを一次情報として扱い、性能superlativeはvendor claimとして明示する。open/deploymentの比較を性能順位へ変換しない。 \cite{src-9156c37fb1c2} \\
\bottomrule
\end{tabularx}
\renewcommand{\arraystretch}{1.0}
\end{center}
\normalsize
\begin{claimboundary}[この比較の境界]
この表は本号で既に検証・参照している一次資料と論文を横断比較の形へ再配置したものである。vendor / project / author が報告した性能値や優位性は独立再現済みの一般則へ変換せず、本文とSource Notesの帰属境界をそのまま維持する。
\end{claimboundary}
